\documentclass[11pt,a4paper]{article}
\usepackage[margin=19mm]{geometry}
\usepackage{fontspec}
\setmainfont{DejaVu Serif}
\setsansfont{DejaVu Sans}
\setmonofont{DejaVu Sans Mono}[Scale=0.85]
\usepackage{amsmath,amssymb}
\usepackage{xcolor}
\usepackage{hyperref}
\hypersetup{colorlinks=true,urlcolor=blue!55!black,pdftitle={Коректний алгоритм двомасштабного моделювання CTMC},pdfauthor={Two-Layer-CTMC}}
\setlength{\parindent}{0pt}
\setlength{\parskip}{5pt}
\newcommand{\block}[1]{\par\medskip{\sffamily\bfseries\large #1}\par\smallskip}
\begin{document}
\raggedright
{\sffamily\bfseries\LARGE Як правильно моделювати\par двомасштабну CTMC\par}
\smallskip
{\sffamily Короткий алгоритм для мережі спільнот • 6 жовтня 2026\par}

\textbf{Рекомендація: один глобальний годинник Gillespie та окремі локальні стани спільнот.} На кожному кроці локальне зараження, одужання й міжспільнотний імпорт конкурують за те, яка подія станеться першою. Час і тип події визначаємо за \textbf{поточним станом}.

\block{1. Що зберігаємо}
Для кожного вузла — стан $S$, $I$ або $R$. Для спільноти $i$ — її граф, розмір $N_i$, кількості $S_i$, $I_i$ та локальні інтенсивності. Для зв'язків спільнот — ваги $W_{ij}$. Також зберігаємо поточний час $t$ і кінцевий час спостереження $H$.

\block{2. Які події можливі зараз}
\begin{itemize}
\setlength{\itemsep}{0pt}\setlength{\parskip}{0pt}
\item \textbf{Локальне зараження вузла $v\in S_i$:}
\[
 a_v=\beta_{\rm in}\sum_{u\in I_i:\,(u,v)\in E_i}w_{uv}.
\]
\item \textbf{Одужання вузла $v\in I_i$:} інтенсивність $\gamma$. Для SI беремо $\gamma=0$.
\item \textbf{Імпорт зі спільноти $i$ у $j$:}
\[
 \lambda_{ij}=\beta_{\rm out}W_{ij}\frac{I_i}{N_i}\frac{S_j}{N_j},\qquad i\ne j.
\]
Це обрана агрегована модель. Після такого імпорту один рівномірно вибраний сприйнятливий вузол у $j$ переходить у стан $I$.
\end{itemize}
Загальна інтенсивність усіх можливих подій:
\[
 A=\sum_i\left(\sum_{v\in S_i}a_v+\gamma I_i\right)+\sum_{i\ne j}\lambda_{ij}.
\]

\block{3. Один крок моделювання}
\begin{enumerate}
\setlength{\itemsep}{1pt}\setlength{\parskip}{0pt}
\item Обчислити або оновити всі змінені інтенсивності та їхню суму $A$.
\item Якщо $A=0$, нових подій немає: завершити моделювання.
\item Згенерувати свіжий $U\sim\mathcal U(0,1)$ і час очікування $\Delta t=-\ln U/A$.
\item Якщо $t+\Delta t>H$, зупинитися на $H$; подію не виконувати.
\item Вибрати \textbf{одну} подію з імовірністю «її інтенсивність / $A$», використавши нове незалежне випадкове число.
\item Перейти до $t\leftarrow t+\Delta t$, виконати подію, записати результат.
\item Повернутися до кроку 1. Імпорт обробляється так само як інша зміна стану.
\end{enumerate}

\newpage
{\sffamily\bfseries\LARGE Імпорт і наступна мікроподія\par}

\block{4. Що відбувається після імпорту}
Імпорт змінює стан вузла та інтенсивності заражень і одужань у спільноті-одержувачі. Оновлюються також усі макроінтенсивності, що залежать від змінених $S_j$ та $I_j$.

\textbf{Після цього генеруємо наступний глобальний крок заново.} У запропонованому алгоритмі немає окремої мікроподії, яку треба пересувати: на кожному кроці вибрано лише одну найближчу подію серед усіх каналів.

Якщо реалізація все ж кешує окрему наступну подію кожної спільноти, то після імпорту в неї дозволено скасувати її старий час, тип і ціль та згенерувати нову локальну подію. Для незмінного стану до наступного зовнішнього втручання:
\[
 \Delta t_i\sim\operatorname{Exp}(\Theta_i),\qquad
 \Theta_i=\sum_{v\in S_i}a_v+\gamma I_i.
\]
Нові числа мають бути незалежними від відкинутого відліку. Кешовані події всіх спільнот і макроподії тоді повинні оброблятися в правильному часовому порядку.

\block{5. Чому повторне генерування правильне}
Експоненційний розподіл не має пам'яті. Якщо $E\sim\operatorname{Exp}(1)$ і подія ще не настала після витрати бюджету $h$, то
\[
 (E-h\mid E>h)\sim\operatorname{Exp}(1).
\]
Тому новий незалежний бюджет дає такий самий закон продовження, як збереження залишку. Це твердження стосується марковських інтенсивностей і подій, що залежать від минулого та поточного стану; майбутня траєкторія для вибору імпорту не використовується.

\block{6. Правила, які забезпечують коректність}
\begin{itemize}
\setlength{\itemsep}{2pt}\setlength{\parskip}{0pt}
\item \textbf{Не відкочувати RNG для повторного використання вже вибраних чисел.} Відкинути подію можна; використані для неї числа вважаються спожитими.
\item \textbf{Не прогнозувати майбутні локальні події для розміщення імпорту в минулому.} Планувальник просувається лише до найближчої реальної події.
\item \textbf{Не підміняти випадкове очікування його середнім $1/A$.} Потрібен саме експоненційний відлік.
\item Один RNG достатній для послідовного глобального алгоритму без відкочувань. Для незалежних рушіїв і прогнозних копій потрібне чітке розділення потоків та їхнього стану.
\item Якщо до $H$ спільноти не досягнуто, фіксуємо цензурування. Значення $H$ не записуємо як фактичний час прибуття.
\end{itemize}

\block{7. Що саме цей алгоритм відтворює точно}
За фіксованих параметрів між подіями алгоритм є точним SSA \textbf{для заданої агрегованої двомасштабної CTMC}. Агрегована інтенсивність $\lambda_{ij}$ та рівномірне розміщення імпорту можуть відрізнятися від передачі справжніми міжспільнотними ребрами. Ця похибка моделі зберігається навіть за точного годинника.

\smallskip
{\footnotesize Основа: прямий метод Gillespie; точність перезапуску після проміжної зміни стану обговорюється в D.~F.~Anderson (2007), \href{https://arxiv.org/html/0708.0370}{розділ VI.1, алгоритм 4}. Цей документ задає рекомендований алгоритм; він не є описом уже виправленого коду.}
\end{document}
